% ============================================================================
% Title: Corrosion in structural materials
% Author: Nathaniel Thomas
% Affiliation:  Texas A&M University, Department of Nuclear Engineering
% Email:        nathaniel@tamu.edu
% Date Created: 3/15/2026
% Version:      1.0
%
% Description:
%   A .tex source file describing the effects of contaminants on molten
%   salt corrosion.
%
% =============================================================================

\section{Effects of Contaminants of Molten Salt Performance}
\label{sec: contaminants}

The presence of contaminants in a molten salt mix causes a series of
undesirable reactions to occur, overall decreasing the performance of the
salt. Among these reactions, one of note is the acceleration of corrosion of
structural metals in direct contact with the salt. Although the effects of
contaminants on salt are not sustained, they can have immediate pronounced
effects on corrosivity \cite{rev23}. \par

\subsection{Oxygen}
\label{sec: oxygen}

Due to their relatively high concentrations and their increased time to remove,
the most concerning contaminants of note are oxides and hydroxide
compounds \cite{rev18}. Since every constituent fluoride salt used in both
FLiBe and FLiNaK is hydroscopic, water from the air will be absorbed into the
salt. Water reacts with the fluorine ion in fluoride salts to form
\ch{HF}, metal
hydroxides, and metal oxides. In salts containing \ch{UF4}, a high enough oxide
concentration was found to cause \ch{UO2} to fall out of solution and collect in
lines, creating hot spots \cite{rev10}, \cite{rev21}. This issue was
minimized by MSRE scientists through the addition of up to \SI{5}{\percent}
\ch{ZrF4} to the salt mixture to
act as a thermodynamically favored oxide buffer \cite{rev21}. The
\ch{HF} formed from the reaction between water and salt can corrode structural
materials such as \ch{Cr} and \ch{Fe} through fluorination. Despite this,
Aia et al., found that high enough concentrations of oxides in molten FLiNaK,
in excess of \SI{1000}{ppm}, can
actually retard the corrosion of 316SS \cite{rev24}. A high oxide concentration
was found to change the corrosion products from \ch{Fe3O4} and \ch{Cr2O3} to
\ch{CrLiO2}, which formed a more stable passivation layer between the salt and
metal \cite{rev24}. \par

Originally, oxygen concentration was measured by reacting salts with an
aggressive fluorinating agent, and measuring the amount of oxygen
released \cite{rev25}. While attempts using \ch{KBrF4} were successful at
measuring up to \SI{98}{\percent} of the oxygen content, aggressive
fluorinating agents present multiple safety issues, as will be discussed in
Section 4 \cite{rev25}, \cite{rev26}. Cyclic voltammetry can be used to
measure a change in the amount of oxygen ions \cite{rev27}; however, since the
exact forms of oxides in the salt are unknown, it is unknown if this method
is sufficient to accurate. Recently, combustion analysis became the standard
for determining oxygen concentration. Combustion analysis works by melting
the salt in a graphite crucible at a high enough temperature such that the
oxygen in the salt will react with the crucible to form \ch{CO}, which can
then be converted to \ch{CO2} and measured via thermal conductivity
detectors \cite{rev26}. \par

\subsection{Sulfur}
\label{sec: sulfur}

Sulfide and sulfate contaminants encourage increased corrosion by altering
redox potential. Sulfur compounds can attack noble metals like nickel and
form \ch{NiS}, which can then be attacked by fluoride ions to form \ch{NiF},
which will dissolve into the salt solution, corroding loop structures
\cite{rev10,rev21}. Although less abundant in constituent salts used
today, sulfur was a major contaminant of consideration for MSRE operations,
due to large amounts of it being present in their \ch{LiF} supply
\cite{rev10}. Sulfates and sulfites also decompose at high
temperatures (\qtyrange{600}{1000}{\celsius}), releasing more oxygen
into the salt \cite{rev29}.  Sulfur contaminants tend to be removed
faster and more completely than oxygen contaminants \cite{rev30}. A standard
method of measuring sulfur content was not found. This is likely due
to the lack of a need for one. Sulfur is completely removed from the
system long before oxygen for all of the purification methods
discussed. \cite{rev27,rev30,rev31,rev32}.\par

\subsection{Structural Metals}
\label{sec: structural metals}

Structural metal contaminants, such as fluorinated \ch{Cr}, \ch{Fe}, and
\ch{Ni}, affect the nuclear cross section of the salt, and can contribute to
loop clogging by deposition to pure metals, and can increase corrosion rates
\cite{rev10,rev29}. Molten fluorides prefer to attack \ch{Cr} over
\ch{Ni} or \ch{Fe} \cite{rev23,rev33,rev34,rev35,rev36,rev37}.
Because \ch{Cr} is the preferred metal to attack, any \ch{Ni} or \ch{Fe} in the
salt as fluorides will shift the redox potential and cause \ch{Cr} to be
further attacked through substitutive fluorination \cite{rev29}. An excess of
\ch{Cr} in the salt can cause it to plate out of solution onto other
metal surfaces, diffusing into those surfaces and weaking the metals resistance
to future corrosion when the redox potential of the salt changes again
\cite{rev29}. Outside of salt corrosion, metal fluorides are often introduced
by the oxidation of structural metals alloy constituents by contact with
strong fluorinating agents \cite{rev21,rev29,rev38}. These
metal fluoride contaminants then dissolve into the molten salt mix to disrupt
its chemical equilibrium. The concentration of metal impurities can be easily
measured through inductively coupled plasma (ICP) spectroscopy \cite{rev26}.
Just like with oxygen, cyclic voltammetry can also be used to determine the
metal concentration relative to a standard but lacks the ability to give
precise values \cite{rev27,rev39} \par
